\pagebreak
\section{Conclusiones}

Se validó experimentalmente que el fraccionamiento del núcleo ferromagnético en chapas delgadas (láminas cuya existencia se comprobó mediante inspección visual en el laboratorio) reduce drásticamente las pérdidas por corrientes de Foucault ($P_f$). Se registró una disminución del $81,2\%$ al pasar de $22,61\text{ W}$ en el núcleo macizo a solo $4,24\text{ W}$ en el laminado. Esto confirma que el aumento de la resistencia transversal, confinada por el aislamiento interlaminar, optimiza la eficiencia energética y reduce el calentamiento térmico crítico bajo excitación variable, mejorando así la eficacia de la laminación en la mitigación de corrientes parásitas.
    
El núcleo macizo mostró severas limitaciones operacionales, alcanzando su corriente límite de seguridad ($3,8\text{ A}$) a una tensión de apenas $40\text{ V}$ debido al efecto de apantallamiento electromagnético provocado por las corrientes de Foucault. Por el contrario, la contención de estas corrientes parásitas en la estructura laminada permitió triplicar la tensión de alimentación (hasta $110\text{ V}$) manteniendo un régimen de corriente menor ($3,0\text{ A}$) y una mayor capacidad de magnetización útil sin comprometer térmicamente los devanados. Con base en esto, se afirma la existencia de restricciones operativas por acoplamiento magnético adverso cuando el núcleo es macizo.
    
La distorsión geométrica en la pantalla del osciloscopio reflejó de forma directa los fenómenos físicos internos de cada núcleo. En el núcleo macizo, el lazo experimentó un ensanchamiento horizontal de carácter ovalado debido a la necesidad de inyectar una mayor corriente de excitación para contrarrestar el campo magnético opuesto inducido por las corrientes parásitas, lo cual corroboró experimentalmente la Ley de Lenz. En el núcleo laminado, al minimizarse este efecto, se visualizó el lazo real del material con sus zonas de saturación magnética claramente delimitadas en los extremos, dando lugar a morfologías notablemente distintas entre ambos ensayos.
    
Se comprobó visualmente mediante el osciloscopio que la mitigación de las corrientes de Foucault en el núcleo laminado expone directamente el carácter no lineal de la curva de magnetización $\vec{B}-\vec{H}$ del material ferromagnético. Esto provoca una marcada distorsión armónica en la forma de onda de la corriente de excitación, la cual exhibe picos pronunciados. Este comportamiento difiere del régimen predominantemente sinusoidal observado en el ensayo del núcleo macizo, donde el efecto parásito dominante enmascara la no linealidad intrínseca del hierro.
    
El uso del circuito integrado pasivo $RC$ resultó ser una herramienta analógica precisa para la transducción de señales. Este circuito permitió obtener una tensión proporcional a la densidad de flujo ($B$) en el capacitor mediante la integración de la fuerza electromotriz (FEM) inducida en el devanado secundario. Asimismo, la caída de tensión en la resistencia shunt proporcionó una señal en fase con la intensidad de campo ($H$), facilitando la reconstrucción exacta del lazo en el modo $X-Y$ del osciloscopio.
    
La aproximación metodológica de despreciar las pérdidas por efecto Joule en el cobre ($P_j \approx 0\text{ W}$) resultó plenamente aceptable para los fines de esta práctica. Dado que las magnitudes de disipación activa en el hierro (tanto por histéresis como por Foucault) representan la componente principal de la potencia activa total registrada por el vatímetro en ambas fases del ensayo, se valida el uso directo de dicha lectura para la caracterización del núcleo ferromagnético.
